\documentclass[12pt,a4paper]{article}
\usepackage[utf8]{inputenc}
\usepackage[T2A]{fontenc}
\usepackage[bulgarian]{babel}
\usepackage{booktabs}
\usepackage{array}
\usepackage{geometry}
\geometry{margin=2.5cm}

\begin{document}

\section*{6.3. Сценарии на експериментите}

\textbf{Сценарий 1: Едри задачи (32 броя)}
\begin{itemize}
    \item 16 тежки задачи: fib(42), fib(43), fib(44), fib(45) $\times$ 4 пъти
    \item 16 леки задачи: fib(35), fib(36), fib(37), fib(38) $\times$ 4 пъти
    \item Нехомогенно разпределение: Процес 0 получава 75\% от тежките задачи
\end{itemize}

\textbf{Сценарий 2: Ситни задачи (128 броя)}
\begin{itemize}
    \item 64 средни + 64 леки задачи (fib 28-32)
    \item Същото нехомогенно разпределение
\end{itemize}

\textit{Защо нехомогенно? Симулира реални сценарии: уеб сървър, научни изчисления.}

\textbf{Сценарий 3: Декомпозиция на Фибоначи (НОВ)}

\textbf{Проблем:} При сценарии 1 и 2 всяка задача fib(n) се изчислява изцяло от един процес.

\textbf{Решение:} fib(45) = fib(44) + fib(43), подзадачите се разпределят между процесите до праг.

\vspace{0.5cm}
\textbf{Избор на оптимален праг за Full Graph (8p, fib 45):}

\begin{center}
\begin{tabular}{|c|c|c|c|c|}
\hline
\textbf{Праг} & \textbf{T1 (s)} & \textbf{T2 (s)} & \textbf{T3 (s)} & \textbf{Средно} \\
\hline
\textbf{28} & 0.408 & 0.371 & 0.370 & \textbf{0.383} \\
\hline
29 & 0.393 & 0.357 & 0.353 & 0.368 \\
\hline
30 & 0.368 & 0.375 & 0.381 & 0.375 \\
\hline
32 & 0.390 & 0.325 & 0.398 & 0.371 \\
\hline
\end{tabular}
\end{center}

\textit{Оптимален праг Full Graph: 28}

\vspace{0.5cm}
\textbf{Избор на оптимален праг за Chain (8p, fib 45):}

\begin{center}
\begin{tabular}{|c|c|c|c|c|c|}
\hline
\textbf{Праг} & \textbf{T1 (s)} & \textbf{T2 (s)} & \textbf{T3 (s)} & \textbf{Средно} & \textbf{Забележка} \\
\hline
28 & \multicolumn{3}{c|}{Блокира при p=1} & -- & Нестабилен \\
\hline
29 & \multicolumn{3}{c|}{Блокира при p=1} & -- & Нестабилен \\
\hline
\textbf{30} & 0.452 & 0.446 & 0.447 & \textbf{0.448} & Работи \\
\hline
32 & 0.490 & 0.460 & 0.500 & 0.483 & По-бавен \\
\hline
\end{tabular}
\end{center}

\textit{Оптимален праг Chain: 30 (минималният стабилен)}

\vspace{0.3cm}
\textbf{Защо Chain блокира при праг < 30?} При единичен процес (p=1) и нисък праг се генерират твърде много подзадачи, което препълва опашката. Full Graph няма този проблем поради по-ефективната комуникация.

\end{document}
